% Part: first-order-logic
% Chapter: syntax-and-semantics
% Section: unique-readability

\documentclass[../../../include/open-logic-section]{subfiles}

\begin{document}

\olfileid[fr]{fol}{syn}{unq}

\olsection{Lecture unique}

\begin{explain}
La manière dont nous avons défini les !!{formula}s garantit que
chacune possède une \emph{lecture unique} : il n'y a essentiellement
qu'une seule façon de la construire selon nos règles de formation,
et une seule façon de l'« interpréter ». Sinon, certaines
!!{formula}s seraient ambiguës et admettraient plusieurs lectures ou
interprétations, ce que nous voulons évidemment éviter. Surtout, sans
cette propriété, la plupart des définitions et des démonstrations qui
suivent ne seraient pas valables.

Pour le comprendre, voyons ce qui se passerait si nous avions donné
des règles de formation ne garantissant pas la lecture unique.
Supposons, par exemple, que nous ayons omis les parenthèses dans la
règle relative aux connecteurs et autorisé ceci :
\begin{quote}
Si $!A$ et $!B$ sont des !!{formula}s, alors $!A \lif !B$ en est une.
\end{quote}
À partir d'une !!{formula} atomique~$!D$, nous pourrions former
$!D \lif !D$, puis, avec~$!D$, obtenir $!D \lif !D \lif !D$. Mais
cette dernière construction peut se faire de deux façons :
\begin{enumerate}
\item Prendre $!A$ pour~$!D$ et $!B$ pour~$!D \lif !D$.
\item Prendre $!A$ pour~$!D \lif !D$ et $!B$ pour~$!D$.
\end{enumerate}
La !!{formula}~$!D \lif !D \lif !D$ admet donc deux lectures.
Elle est de la forme $!B \lif !C$ avec $!B$ égal à~$!D$ et $!C$
égal à~$!D \lif !D$ ; mais elle est \emph{aussi} de la forme
$!B \lif !C$
avec $!B$ égal à~$!D \lif !D$ et $!C$ égal à~$!D$.

Dans ce cas, nos définitions ne fonctionneraient pas toujours. Par
exemple, pour définir l'!!{main operator} d'une !!{formula}, nous
disons que, dans une !!{formula} de la forme $!B \lif !C$, il s'agit
de l'occurrence indiquée de~$\lif$. Or, si $!D \lif !D \lif !D$
correspond à $!B \lif !C$ des deux façons ci-dessus, la première
occurrence de~$\lif$ est l'!!{main operator} dans un cas, et la
seconde dans l'autre. Nous voulons pourtant que l'!!{main operator}
soit une \emph{fonction} de la !!{formula} : chaque !!{formula} non
atomique doit posséder exactement une occurrence de son !!{main operator}.
\end{explain}

\begin{lem}
Dans !!a{formula}~$!A$, le nombre de parenthèses ouvrantes est égal
au nombre de parenthèses fermantes.
\end{lem}

\begin{proof}
Nous raisonnons par induction sur la construction de~$!A$. Il faut
d'abord vérifier la propriété pour toutes les !!{formula}s atomiques
(cas de base), puis montrer que, si elle vaut pour des !!{formula}s
données, elle vaut aussi pour celles que l'on construit à partir
d'elles (étape d'induction).

Notons $l(!A)$ le nombre de parenthèses ouvrantes et $r(!A)$ celui
des parenthèses fermantes dans~$!A$ ; définissons de même $l(t)$ et
$r(t)$ pour un terme~$t$.

\begin{prob}
    Montrer que $l(t) = r(t)$ pour tout terme~$t$.
\end{prob}

\begin{enumerate}
\tagitem{prvFalse}{\indcase{!A}{\lfalse}{$\indfrm$ contient $0$
    parenthèse ouvrante et $0$ parenthèse fermante.}}{}

\tagitem{prvTrue}{\indcase{!A}{\ltrue}{$\indfrm$ contient $0$
    parenthèse ouvrante et $0$ parenthèse fermante.}}{}

\item \indcase{!A}{\Atom{R}{t_1,\dots,t_n}}{$l(\indfrm) = 1 + l(t_1) +
  \dots + l(t_n) = 1 + r(t_1) + \dots + r(t_n) = r(\indfrm)$. Nous
  utilisons ici le fait, laissé en exercice, que $l(t) = r(t)$ pour
  tout terme~$t$.}

\item \indcase{!A}{\eq[t_1][t_2]}{$l(\indfrm) = l(t_1) + l(t_2) =
  r(t_1) + r(t_2) = r(\indfrm)$.}

\tagitem{prvNot}{\indcase{!A}{\lnot !B}{Par hypothèse d'induction,
    $l(!B) = r(!B)$. Donc $l(\indfrm) = l(!B) = r(!B) =
    r(\indfrm)$.}}{}

\item \indcase{!A}{(!B \ast !C)}{Par hypothèse d'induction, $l(!B) =
  r(!B)$ et $l(!C) = r(!C)$. Donc $l(\indfrm) = 1 + l(!B) + l(!C) =
  1 + r(!B) + r(!C) = r(\indfrm)$.}

\tagitem{prvAll}{\indcase{!A}{\lforall[x][!B]}{Par hypothèse
    d'induction, $l(!B) = r(!B)$. Donc $l(\indfrm) = l(!B) = r(!B) =
    r(\indfrm)$.}}{}

% Afficher le cas existentiel si prvEx et non prvAll.
\tagitem{prvAll,notprvEx}{}{\indcase{!A}{\lexists[x][!B]}{Par
    hypothèse d'induction, $l(!B) = r(!B)$. Donc $l(\indfrm) = l(!B)
    = r(!B) = r(\indfrm)$.}}

% Dire « de même » si prvEx et prvAll.
\tagitem{notprvAll,notprvEx}{}{\indcase{!A}{\lexists[x][!B]}{De même.}}
\end{enumerate}
\end{proof}

\begin{defn}[Préfixe propre]
Une chaîne de symboles~$!B$ est un \emph{préfixe propre} d'une chaîne
de symboles~$!A$ si la concaténation de~$!B$ avec une chaîne non vide
donne~$!A$.
\end{defn}

\begin{lem}\ollabel{lem:no-prefix}
Si $!A$ est !!a{formula} et $!B$ un préfixe propre de~$!A$, alors
$!B$ n'est pas !!a{formula}.
\end{lem}

\begin{proof}
Exercice.
\end{proof}

\begin{prob}
Démontrer le \olref[fol][syn][unq]{lem:no-prefix}.
\end{prob}

\begin{prop}
\ollabel{prop:unique-atomic}
Si $!A$ est !!a{formula} atomique, elle vérifie une et une seule des
conditions suivantes.
\begin{enumerate}
\tagitem{prvFalse}{$!A \ident \lfalse$.}{}
\tagitem{prvTrue}{$!A \ident \ltrue$.}{}
\item $!A \ident \Atom{R}{t_1,\dots,t_n}$, où $R$ est un !!{predicate}
  à $n$ places, $t_1$,~\dots, $t_n$ sont des termes ; $R$ et chacun
  des termes $t_1$,~\dots, $t_n$ sont déterminés de façon unique.
\item $!A \ident \eq[t_1][t_2]$, où $t_1$ et $t_2$ sont des termes
  déterminés de façon unique.
\end{enumerate}
\end{prop}

\begin{proof}
Exercice.
\end{proof}

\begin{prob}
Démontrer la \olref[fol][syn][unq]{prop:unique-atomic}. Indication :
formuler et démontrer pour les termes une version du
\olref[fol][syn][unq]{lem:no-prefix}.
\end{prob}

\begin{prop}[Lecture unique]
Toute !!{formula} vérifie une et une seule des conditions suivantes.
\begin{enumerate}
\item $!A$ est atomique.

\tagitem{prvNot}{$!A$ est de la forme $\lnot !B$.}{}

\tagitem{prvAnd}{$!A$ est de la forme $(!B \land !C)$.}{}

\tagitem{prvOr}{$!A$ est de la forme $(!B \lor !C)$.}{}

\tagitem{prvIf}{$!A$ est de la forme $(!B \lif !C)$.}{}

\tagitem{prvIff}{$!A$ est de la forme $(!B \liff !C)$.}{}

\tagitem{prvAll}{$!A$ est de la forme $\lforall[x][!B]$.}{}

\tagitem{prvEx}{$!A$ est de la forme $\lexists[x][!B]$.}{}
\end{enumerate}
De plus, dans chaque cas, $!B$, ou bien $!B$ et $!C$, sont déterminées
de façon unique. Ainsi, il n'existe pas deux paires distinctes $!B$,
$!C$ et $!B'$, $!C'$ telles que $!A$ soit à la fois de la forme
\iftag{prvIf}{$(!B \lif !C)$ et $(!B' \lif !C')$.}{
    \iftag{prvOr}{$(!B \lor !C)$ et $(!B' \lor !C')$.}{
      \iftag{prvAnd}{$(!B \land !C)$ et $(!B' \land !C')$.}{}}}
\end{prop}

\begin{proof}
Les règles de formation imposent à toute !!{formula} non atomique de
commencer par une parenthèse ouvrante~(, \iftag{prvNot}{$\lnot$,}{}
ou par un quantificateur. Toute !!{formula} qui commence par l'un des
symboles suivants est, en revanche, nécessairement atomique :
!!a{predicate}, !!a{function}, !!a{constant}, !!a{variable}%
\iftag{prvFalse}{, $\lfalse$}{}\iftag{prvTrue}{, $\ltrue$}{}.

Il nous reste donc seulement à montrer que si $!A$ est de la forme
$(!B \ast !C)$ et aussi de la forme $(!B' \mathbin{\ast'} !C')$,
alors $!B \ident !B'$, $!C \ident !C'$ et $\ast = {\ast'}$.

Supposons donc que $!A \ident (!B \ast !C)$ et
$!A \ident (!B' \mathbin{\ast'} !C')$. Soit $!B \ident !B'$, soit
ce n'est pas le cas. Dans le premier cas, $\ast = {\ast'}$ et
$!C \ident !C'$ : ce sont alors des sous-chaînes de~$!A$ qui
commencent au même endroit et ont la même longueur. Dans le second
cas, $!B \not\ident !B'$. Puisque $!B$ et $!B'$ sont des sous-chaînes
de~$!A$ commençant au même endroit, l'une est un préfixe propre de
l'autre. Or cela contredit \olref{lem:no-prefix}.
\end{proof}

\end{document}
